Short-read metagenomic next-generation sequencing (mNGS) pipelines process trimmed, ambiguous or locally perturbed reads. Because representation compression precedes classification or database lookup, downstream accuracy alone does not reveal which sequence attributes were retained. We present a representation-diagnostics framework with a dense audit from 50 to 75 bp and anchors at 100, 125 and 150 bp. CK4P-MSP combines a reverse-complement canonical 4-mer block (K), a global property block (P) and multi-scale property pooling (MSP) in a 222-dimensional block-normalized representation. A seven-group matched ablation assigned P primarily to global stability, MSP to grouped local-change readability and K to composition-linked retrieval. In the prespecified 69--150 bp grid, CK4P-MSP had mean paired cosine 0.989, mean drift 0.129 and grouped local-versus-noise macro-F1 of 0.979. A shared-template sweep over every integer length from 50 to 75 bp retained the same operating profile. Across six CAMI target-versus-background tasks, a probe trained only on 100-bp clean reads achieved mean shifted macro-F1 of 0.793 for CK4P-MSP, comparable to CK4/CK5 and 0.010--0.012 above PseKNC/PseEIIP; train-fit coordinate scaling exposed a separate N-mask failure mode. PseKNC achieved lower drift, whereas CK4P-MSP retained stronger grouped local-change readability. Full-position encodings remained higher-dimensional upper bounds for fine positional information, and compressed high-k controls were less stable at the same feature budget. CK4P-MSP therefore occupies an attributable stability--readability trade-off rather than minimizing drift or maximizing predictive accuracy alone. Compact biochemical and coarse positional summaries can provide interpretable representation-level audit coordinates alongside database-linked exact matching and alignment.
